\chapter{Longitud, distancia y geodésicas}
\label{cap:longitud-distancia}

Una métrica convierte velocidades coordenadas en longitudes geométricas. En los modelos anteriores el mismo principio aparece con expresiones diferentes. Este capítulo muestra cómo pasar de una curva parametrizada a una longitud, cómo reconocer una geodésica y cómo verificar cálculos con ejemplos.

\section{Longitud de una curva}

Sea \(\gamma(t)=(x(t),y(t))\), con \(a\leq t\leq b\), una curva a trozos suave en el semiplano superior. A partir de la métrica \cref{eq:metrica-semiplano}, su longitud hiperbólica es
\begin{equation}
  L_{\mathbb H}(\gamma)=\int_a^b
  \frac{\sqrt{x'(t)^2+y'(t)^2}}{y(t)}\,dt.
  \label{eq:longitud-semiplano}
\end{equation}
La expresión no cambia si se reparametriza la curva sin invertir su orientación. Para el disco se reemplaza el factor \(1/y\) por \(2/(1-|z|^2)\).

\section{Un cálculo vertical}

Considera el segmento vertical entre \(z_1=iy_1\) y \(z_2=iy_2\), con \(y_1,y_2>0\). Tomando \(x(t)=0\) y \(y(t)=t\) en \cref{eq:longitud-semiplano}, se obtiene
\begin{equation}
  d_{\mathbb H}(iy_1,iy_2)=\left|\int_{y_1}^{y_2}\frac{dt}{t}\right|
  =\left|\log\frac{y_2}{y_1}\right|.
  \label{eq:distancia-vertical}
\end{equation}
Por ejemplo, la distancia entre \(i\) y \(3i\) es \(\log 3\). El segmento vertical es geodésico, por lo que realiza la distancia mínima.

\begin{example}[Distancia desde el centro del disco]
En el disco, la distancia del origen al punto real \(r\in[0,1)\) es
\[
  d_{\mathbb D}(0,r)=\int_0^r\frac{2\,dt}{1-t^2}
  =\log\frac{1+r}{1-r}=2\operatorname{artanh}(r).
\]
Para \(r=1/2\), esta distancia vale \(\log 3\), igual que la distancia entre \(i\) y \(3i\) calculada en \cref{eq:distancia-vertical}. Una isometría entre los modelos explica la coincidencia.
\end{example}

\section{La ecuación de las geodésicas}

Una geodésica parametrizada a velocidad constante satisface la ecuación de Euler--Lagrange asociada a la energía de la métrica. Para \(ds^2=(dx^2+dy^2)/y^2\), la energía lagrangiana es
\[
  \mathcal L=\frac{x'^2+y'^2}{2y^2}.
\]
Como \(x\) no aparece explícitamente, el momento \(x'/y^2\) es constante. Al combinar esta primera integral con la condición de velocidad constante, se obtienen rectas verticales y circunferencias con centro en el eje real. La comprobación geométrica y el modelo del disco se estudian en los \cref{cap:semiplano-superior,cap:disco-poincare}.

\begin{problem}
Usa \cref{eq:distancia-vertical} para calcular la distancia entre \(2i\) y \(5i\). Explica por qué intercambiar los extremos no cambia el resultado.
\end{problem}

\begin{remark}
La distancia entre puntos es el ínfimo de las longitudes de curvas que los unen. La fórmula de longitud asigna una longitud a cada curva, pero no afirma por sí sola que cualquier curva sea la más corta.
\end{remark}
